% Method: technical core.

\subsection{Problem Formulation}
\label{sec:problem}

We consider an agent $\pi_\theta$ that interacts with an environment
over discrete steps $t = 1, 2, \dots, T$. At each step the agent
observes a state $s_t$, consults a context--memory layer
$\mathcal{C}_t$, produces an action $a_t$ (a natural-language response,
a tool call, or a context operation), and updates both the environment
and the layer. A \emph{context--memory layer} is a pair
$(\mathcal{K}, \mathcal{Q})$ where $\mathcal{K}$ is a set of keyed
entries and $\mathcal{Q}$ is a set of query operators. Entries may be
short-lived working context, long-term memory, tool metadata, retrieved
knowledge, artifacts, tasks, or audit evidence; the layer manages them
through one uniform interface. In a flat history, $\mathcal{K}$ is a
totally ordered sequence and $\mathcal{Q}$ is concatenation. In vector
RAG, $\mathcal{K}$ is embedded and $\mathcal{Q}$ is top-$k$ similarity.
We take $\mathcal{K}$ to be a forest of hierarchical paths and let
$\mathcal{Q}$ include exact lookup, prefix scan, glob match, and
semantic search, each parameterised by a disclosure level.

The agent's objective is a task-defined reward $R$ under a total
token budget $B$. Let $\tau_t$ denote the tokens consumed at step
$t$ by the materialised context. Our design targets the Pareto
frontier of $R$ against $\sum_t \tau_t$: an agent should be able to
achieve comparable $R$ at substantially smaller $\sum_t \tau_t$ than
a flat-history baseline.

\subsection{Overview}
\label{sec:overview}

\textsc{Structure} implements prompt construction as a materialisation
contract over two stores. Durable entries form the candidate set for
long-context disclosure; typed events form the candidate set for
short-context visibility. The system has three implementation layers. An
\emph{addressing layer} maps every managed long-context item to a path
such as \texttt{/tools/web\_search/schema} or
\texttt{/knowledge/python\_guide/chapter\_3}. A \emph{disclosure layer}
exposes each entry at three levels via a uniform \texttt{disclose}
method. An \emph{audit layer} emits every mutation as an immutable
event on a log that backs both real-time streaming, short-context
compression, and deterministic replay. The three layers are
orthogonal: an entry has a path, can be disclosed, and is audited
regardless of which concrete store holds it.

% \begin{figure}[t]
%     \centering
%     \includegraphics[width=\columnwidth]{figures/overview.pdf}
%     \caption{Overview of the \textsc{Structure} context layer. The
%     agent issues typed queries against the addressing layer
%     (glob/prefix/semantic), receives entries at an explicitly chosen
%     disclosure level, and mutates state through operations that are
%     published to the audit log.}
%     \label{fig:overview}
% \end{figure}

\subsection{Context Materialisation Contract}
\label{sec:long_short_context}

The central design choice in \textsc{Structure} is not that agents need
retrieval, history, tools, or memory; those are already standard. The
contribution is a reproducible rule for combining them into the prompt.

\paragraph{Long-context candidates.}
Any state that may be useful beyond the current plan--act cycle is
stored as an addressable entry. This includes retrieved documents
(\texttt{/knowledge/...}), durable conversation facts
(\texttt{/memory/...}), tool schemas and usage notes
(\texttt{/tools/...}), generated files (\texttt{/artifacts/...}), task
state (\texttt{/tasks/...}), and run observations
(\texttt{/runs/<id>/observations/...}). The source does not determine
the interface: a RAG chunk, a history summary, and a tool schema are
all queried by path, prefix, glob, or semantic search and are all
materialised through \textsc{glance}, \textsc{overview}, or
\textsc{detail}. These entries are candidates for the prompt, not
automatic prompt content.

\paragraph{Short-context visibility.}
Events are the unit of short context. Each user message, assistant
reply, tool call, tool result, retrieval, source rating, and artifact
mutation is written to the append-only log with
a type, timestamp, sequence number, and payload. Before every model
step, a materialiser constructs the visible short context in four
passes: (1) keep pinned anchors such as the opening task, explicit user
constraints, final decisions, and artifact mutations; (2) keep recent
plan--act events whose type-specific TTL has not expired; (3) drop
noise events such as token deltas, heartbeats, and superseded
intermediate traces; and (4) protect a recency floor of the most
recent $k$ events regardless of type. The unpinned portion of the
result is bounded by the TTL policy, while pinned anchors accumulate
for the life of the run, so the working context grows with durable
milestones rather than with the full trajectory.

\paragraph{Prompt assembly.}
At step $t$, the prompt is assembled from
$M_t = \mathrm{gc}(\mathcal{E}, t, \Pi) \cup
\mathrm{select}(\mathcal{K}, q_t, \ell_t)$, where
$\mathrm{gc}$ returns the compressed event slice and
$\mathrm{select}$ returns long-context entries from the path space at
disclosure level $\ell_t$. Thus the paper's testable object is the
materialiser: which events remain visible, which path entries are
selected, and at what disclosure level.
Expired short-context events are not lost: because they remain in the
log, later executors can summarise, pin, or replay them, and
evaluators can reconstruct the exact prompt from the event policy and
path selections.

\subsection{Path-Addressable Context}
\label{sec:component_a}

Entries are stored in a relational table with a canonical \texttt{path}
column indexed by a B-tree. The path is a POSIX-style string whose
semantics are defined by convention rather than by the store: tools
live under \texttt{/tools}, knowledge under \texttt{/knowledge}, and
workspace artefacts under \texttt{/workspaces/<id>/\dots}. Three
operators answer $\mathcal{Q}$:
\begin{itemize}
    \item \textbf{Exact} --- \texttt{get(path)} returns at most one
    entry.
    \item \textbf{Prefix} --- \texttt{list(prefix, recursive)} returns
    the set of entries whose path begins with \texttt{prefix}; the
    underlying query is a single indexed range scan.
    \item \textbf{Glob} --- \texttt{glob(pattern)} supports \texttt{*}
    and \texttt{**} wildcards. We first compute the longest non-glob
    prefix of the pattern and issue a prefix scan, then filter in
    application code against the compiled regex. This keeps the
    database work linear in the size of the result set, not of the
    workspace.
\end{itemize}
A parallel set of HNSW vector indices supports semantic retrieval at
four dimensionalities (384, 768, 1024, 1536), letting callers trade
cost and latency against embedding quality without duplicating
storage.

Each entry carries a \emph{scope} --- \textsc{user},
\textsc{workspace}, or \textsc{global} --- that governs visibility.
A single query returns the union across scopes through a SQL
disjunction (Listing~\ref{lst:scope-load}):
shared tool schemas live at \textsc{global}, team knowledge at
\textsc{workspace}, and private notes at \textsc{user}, without the
agent having to issue three separate queries.

\begin{figure}[t]
\small
\begin{verbatim}
stmt = select(Context).where(
  or_(
    and_(Context.source_id == workspace_id,
         Context.scope == WORKSPACE),
    and_(Context.user_id  == owner_id,
         Context.scope == USER),
    Context.scope == GLOBAL,
  ))
\end{verbatim}
\caption{Multi-scope load. Three visibility levels are resolved in a
single indexed query and presented to the agent as one list.}
\label{lst:scope-load}
\end{figure}

\subsection{Multi-Level Disclosure}
\label{sec:component_b}

The disclosure layer converts a stored entry into one of three
projections. We define the projection as a method on the entry
itself, not as a separate summarisation pipeline, so the three levels
are always consistent with the stored state:
\begin{itemize}
    \item \textsc{glance} --- a short summary ($\leq 50$\,B by
    default) intended to answer \emph{``does this matter to me now?''}
    An agent can scan hundreds of glances in a single prompt.
    \item \textsc{overview} --- glance plus structured metadata
    (tags, type, scope, timestamps), adequate for routing and
    decision-making without committing to full content.
    \item \textsc{detail} --- full content plus unstructured metadata.
    Large payloads are stored in S3 and resolved lazily.
\end{itemize}
Each query in $\mathcal{Q}$ accepts the level as an argument, so the
agent's policy is \emph{jointly} a choice of address pattern and
disclosure depth. This is the central mechanism by which the layer
controls $\sum_t \tau_t$: the agent reads glances until it needs to
commit, at which point it escalates a narrow subset to \textsc{detail}.

\begin{figure}[t]
\small
\begin{verbatim}
def disclose(self, level="overview"):
  out = {"path": self.path,
         "glance": self.glance or self.content[:50]}
  if level == "glance":
    return out
  if self.tags: out["tags"] = self.tags
  if level == "overview":
    return out
  out["content"] = self.content
  out["meta"]    = self.meta or {}
  return out
\end{verbatim}
\caption{Disclosure as a method on the entry. The three levels are
monotone: each extends the previous one.}
\label{lst:disclose}
\end{figure}

\subsection{Source Evaluation via Agent Ratings}
\label{sec:component_rating}

Not all context entries are equally useful at inference time: a tool
schema may be out of date, a retrieved document may be off-topic, and an
agent-generated note from an earlier step may be wrong. Existing memory
systems treat these cases implicitly --- either the entry is retrieved or
it is not. We make the judgement explicit: after consulting an entry,
the agent can call a \texttt{rate\_context(path, rating, comment)}
operation with a bounded rating $r \in [-1, 1]$ (positive = helpful,
negative = misleading, zero = neutral). Each call emits a
\textsc{context.rated} event and updates a denormalised pair
$(\text{rating\_sum}, \text{rating\_count})$ on the entry, from which the
mean $\bar r$ is derived.

The write path above --- the rating operation, the
\textsc{context.rated} event, and the denormalised aggregate --- is
what the released platform implements. The read path is specified as
part of the contract and consumes $\bar r$ in two ways. At the filter
stage, a threshold $\tau$ discards entries the agent or prior runs
have consistently judged harmful. At the ranking stage, a convex
combination
$s = \alpha \cdot s_\text{sim} + (1 - \alpha) \cdot \bar r$ blends
semantic similarity $s_\text{sim}$ with the rating mean, where
$\alpha \in [0, 1]$ is a per-query hyperparameter defaulting to
$\alpha = 0.7$; the rating ablation (\S\ref{sec:results}) evaluates
this read path against a no-rating arm.
Because ratings are event-sourced, the per-step state is
reproducible: replaying events up to step $t$ reconstructs the exact
$\bar r$ the agent saw at that step. The mechanism is closest in spirit
to the reflection tokens of Self-RAG \citep{asai2024selfrag} and the
verbal feedback of Reflexion \citep{shinn2023reflexion}, but applied to
the \emph{source} rather than to the agent's own output and stored in
a structured aggregate rather than in free-form text.

\subsection{Event-Level Garbage Collection}
\label{sec:component_gc}

Even with multi-level disclosure, an agent's working context grows with
each step: every tool call, every intermediate thought, every retrieval
result produces events. Without pruning, the prompt either exceeds the
window or pushes high-value turns into the middle-of-context
\citep{liu2024lostmiddle}. Blind truncation by age loses artifacts;
free-form summarisation hides structure.

We introduce an in-loop garbage collector parameterised by an
\emph{event time-to-live} (TTL) per event type. The collector is a pure
function $\mathrm{gc}(\mathcal{E}, t, \Pi) \to \mathcal{E}_{\text{live}}$
that takes a sequence of events $\mathcal{E}$, the current step $t$, and
a policy $\Pi: \mathcal{T} \to (\text{ttl}, \text{pin})$
defined over event types $\mathcal{T}$, and returns the subset still
visible. An event $e$ of type $\tau$ ages by the events that arrive
after it: each newer event contributes a decay increment (default
$1$), and type-specific decay rules let designated successors age an
event faster --- for example, a \texttt{tool.result} adds $3$ to the
decay of the preceding \texttt{tool.call}, so answered calls expire
promptly. $e$ is kept iff $\Pi(\tau).\text{pin}$ is true, its
accumulated decay is at most $\Pi(\tau).\text{ttl}$, or it falls
within a recency floor that keeps the most recent $k$ events
unconditionally ($k{=}100$ per run, $k{=}500$ per workspace by
default).
Crucially, events are \emph{not} deleted from storage; only their
visibility to the agent is gated, so the audit log and offline replay
remain exact.

The default policy (Appendix, Table~\ref{tab:ttl}) pins user and
assistant messages, run lifecycle milestones, artifact mutations, task
completions, and system errors for the life of the run;
gives tool calls and results a short event-TTL of roughly the width of
a single plan--act cycle; and assigns zero TTLs to high-frequency
noise events (tokens, heartbeats) so they never enter the prompt. When
a span has expired but remains semantically useful, an executor can
promote it: write a structured summary as a long-context entry under
the run's path (via the ordinary \texttt{create} operator) and let
disclosure bring it back at \textsc{glance} level. The collector
itself performs no summarisation.
Executors may override any entry of $\Pi$ per run. In the
pathological case where all events are pinned the collector degenerates
to the flat-history baseline; in the other extreme all but user and
artifact events expire immediately, approximating a ReAct-style
stateless loop \citep{yao2023react}. Everything between those endpoints
is configuration.

\paragraph{Comparison with concurrent work.}
AgentDiet \citep{xu2025agentdiet} invokes a reflection LLM call past a
step threshold to condense earlier trajectory content, reporting
$40$--$60\%$ input-token reduction on SWE-agent-style tasks.
Memory-as-Action \citep{zhang2025memact} lifts memory editing into the
agent's action space as a single \texttt{prune\_context(summary, ids)}
tool and trains the policy end-to-end with a trajectory-segmented RL
objective (DCPO), reaching $51\%$ context reduction on Qwen2.5-14B.
Context-Folding \citep{sun2025contextfolding} and AgentFold
\citep{ye2025agentfold} fold entire sub-trajectories into summary
nodes, using RL and SFT respectively. MemAgent
\citep{yu2025memagent} optimises a segment-reading and memory-overwrite
workflow for extreme long-context QA, demonstrating that learned memory
policies can extend effective context length when retraining is
available. All five operate on the
\emph{time} axis of a single trajectory and \emph{destructively} rewrite
it, either paying a per-step LLM call (AgentDiet, MemAct) or a
model-specific training cost (MemAct, Context-Folding, AgentFold,
MemAgent).
Our collector occupies a different point in the design space: it
operates on the \emph{address} axis of a persistent workspace via a
\emph{declarative, zero-compute} function whose output is a visibility
filter, not a mutation of storage. Because the event log is untouched,
every $\mathrm{gc}$ result at step $t$ is byte-exactly reproducible
from the configuration $\Pi$ alone, and any expired event remains
available to a later disclosure or replay. We view the learned-policy
methods above as candidates to be layered \emph{on top of} our
substrate --- a trained pruner or memory-overwrite policy can express
its pin, expire, and summarise decisions as first-class events on the
same log, combining policy quality with auditability.

\subsection{Event-Sourced Audit Layer}
\label{sec:component_c}

Every creation, update, and deletion of an entry produces an event
whose payload describes the mutation. Events are assigned a
per-run and per-workspace sequence number and written to
PostgreSQL; the same events are streamed to Redis for real-time
delivery to connected clients. An agent's context trajectory is
therefore a totally ordered sequence that can be replayed
deterministically --- useful for debugging, for offline evaluation,
and for benchmark reproducibility. Sequence numbers are maintained by
an in-memory counter seeded from a single
\texttt{SELECT MAX(sequence)} at cold start; subsequent events pay
only an increment.

\subsection{Runtime Pipeline}
\label{sec:runtime}

A typical step proceeds as follows. (1) The agent issues a
\texttt{glance(/tools)} query and receives, for example, fifteen
glances in roughly 3\,KB. (2) Conditioned on the glances, it issues a
\texttt{detail(\allowbreak/tools/web\_search)} for one entry and
receives its full schema. (3) It emits a tool call; the resulting
output is persisted as a new context entry under
\texttt{/workspaces/<id>/runs/<run\_id>/observations} and the event is
broadcast. (4) A subsequent reasoning step reads only the glance of
the observation unless the agent chooses to revisit the detail. Across
a long trajectory, the ratio of detail reads to glance reads
determines token cost.

\subsection{Implementation Notes}

We implement \textsc{Structure} as a Python service with an async
FastAPI surface, a PostgreSQL-backed store (pgvector for embeddings),
and a Redis Streams event bus. Ten first-class tools expose the
operators above to the agent: \texttt{glance}, \texttt{list},
\texttt{tree}, \texttt{glob}, \texttt{read}, \texttt{search},
\texttt{create}, \texttt{update}, \texttt{delete}, and \texttt{rate}
for source evaluation. A per-worker in-memory cache,
invalidated through a Redis dirty flag, reduces warm path lookups
from roughly $500$\,ms to under $10$\,ms.

To keep the paper claim aligned with the released implementation, the
repository contains a canonical core manifest at
\path{core/structure_core.json}. The manifest names the paper-level
primitives above (path addressing, disclosure, audit events, source
evaluation, event-level GC, and reproducible evidence) and maps each to
the three operational surfaces used in the artifact: the service-backed
web application, the local desktop workbench, and the scriptable
CLI/TUI. The web application imports the manifest as a read-only
contract while preserving its API-backed runtime; the desktop app and
CLI/TUI consume the same manifest through the Rust local-core crate.
The manifest additionally contains a capability matrix: each capability
is linked to a paper primitive and records, for every surface, whether
the implementation is native to that surface, service-backed, or
read-only, together with the concrete entrypoint and evidence used to
verify it. This makes the intended asymmetry explicit: CLI/TUI,
desktop, and web do not expose identical workflows, but they preserve
the same conceptual contracts for context addressing, event replay, and
reproducible evidence. The same manifest also names the benchmark
report evidence schema (\texttt{benchmark-report-v1}), keeping
service-backed benchmark adapters and embedded local benchmark runs
aligned on cost, per-case, diagnostic, and evidence fields. For the
local artifact, the Rust runtime can export a single evidence bundle
(\texttt{local-evidence-bundle-v1}) containing the manifest snapshot,
surface-parity report, workspace replay, per-run evidence summaries,
knowledge-source metadata, artifact records, and recent benchmark
reports; this bundle is exposed consistently through the CLI/TUI and
desktop app. Individual benchmark report files can also be re-linked
to their underlying run evidence through a
\texttt{local-benchmark-evidence-v1} export, so the reported score,
case diagnostics, event stream, tool calls, and generated artifacts
remain inspectable from the terminal UI, the scriptable CLI, and the
desktop app.
